\documentclass[10pt,a4paper]{article}
\usepackage[spanish,es-nodecimaldot]{babel}
\usepackage[utf8]{inputenc}
\usepackage[T1]{fontenc}
\usepackage{lmodern}
\usepackage[top=1.1cm,bottom=1.1cm,left=1.5cm,right=1.5cm]{geometry}
\usepackage{amsmath,amssymb}
\usepackage[table]{xcolor}
\usepackage{booktabs,array}
\usepackage{enumitem}

\setlength{\parindent}{0pt}
\setlength{\parskip}{2pt}
\pagestyle{empty}

\AtBeginDocument{%
  \setlength{\abovedisplayskip}{4pt}%
  \setlength{\belowdisplayskip}{4pt}%
  \setlength{\abovedisplayshortskip}{2pt}%
  \setlength{\belowdisplayshortskip}{2pt}%
}

\begin{document}

\begin{center}
{\large\bfseries Metodo de la varianza residual}
\end{center}
\vspace{-2pt}

Para evitar el sobreajuste, se calcula la \textbf{varianza residual}:
\[
\sigma_{\text{res}}^2 = \frac{\displaystyle\sum_{i=1}^{n} (y_i-\hat{y}_i)^2}{n-(p+1)}
\]

Donde:
\begin{itemize}[leftmargin=16pt,itemsep=1pt,topsep=1pt,parsep=0pt]
  \item $n$ es la cantidad de puntos de la muestra.
  \item $p$ es el grado del polinomio.
  \item $p+1$ es la cantidad de coeficientes estimados.
  \item $n-p-1$ representa los \textbf{grados de libertad del modelo}.
  \item $y_i - \hat{y}_i$ es el residuo, es decir, la diferencia entre el valor observado y el valor estimado.
\end{itemize}

\textbf{Se comienza por :}
\begin{itemize}[leftmargin=16pt,itemsep=0pt,topsep=1pt,parsep=0pt]
  \item una recta, grado 1
  \item se sigue parabola, grado 2,
  \item luego polinomio, grado 3
  \item polinomio de grados superiores si fuera necesario
\end{itemize}
Para cada uno se calcula la varianza residual y se compara los resultados.
La suma de cuadrados de error vista en el numerador de la formula tiende a disminuir porque la curva se adapta mejor a los puntos
Tambien disminyen los grados de libertad, porque se estiman mas coeficientes

\textbf{Para la seleccion:}
Se selecciona el grado en el q se observa la mayor caida. Despues de este punto aumentar el grado puede ser innecesario ya q puede aumenta el riesgo de sobreajuste

\textbf{Ej:}

Dada una muestra de $n=10$ mediciones industriales, se evalúa el ajuste de polinomios de grados sucesivos obteniéndose los siguientes resultados reales:

\begin{center}
{\small
\begin{tabular}{c c c c p{3.6cm}}
\toprule
\textbf{Grado (p)} & \textbf{SCE Muestral} & \textbf{GL ($n-p-1$)} & \textbf{Var. Residual ($s^2$)} & \textbf{¿Se justifica?} \\
\midrule
1 (Lineal) & 57.655 & 8 & \textbf{7.207} & Modelo Base \\
2 (Cuadrático) & 53.059 & 7 & \textbf{7.580} & No, la varianza aumentó \\
3 (Cúbico) & 6.124 & 6 & \textbf{1.021} & \textbf{Sí, salto decreciente masivo} \\
4 (Cuártico) & 4.959 & 5 & \textbf{0.992} & No, reducción insignificante \\
\bottomrule
\end{tabular}
}
\end{center}

\end{document}
